% ============================================================================
% DRAFT manuscript additions - new measured results (2026-08-09).
% Factual/table content only; prose kept minimal for the author to humanize
% (Turnitin). 100% ASCII, no em dashes. Drop into frontier.tex where noted.
% ============================================================================

% ---------------------------------------------------------------------------
% TABLE: measured on-device energy per inference (replaces the 50 mA estimate)
% Reviewer Major #8. Classic ESP32-D0WD-V3, 240 MHz, WiFi off, measured with an
% INA219 in series with the 5 V board input; an external ESP8266 logs bus voltage
% and current, and a GPIO marker driven HIGH during each Invoke() separates
% active-inference power from the idle baseline (254 mW / 50.6 mA at 5 V).
% ---------------------------------------------------------------------------
\begin{table}[t]
\centering
\caption{Measured on-device energy per inference on the classic ESP32
(240 MHz, WiFi off). Active power is the mean bus power while the marker is high;
the idle baseline is 254 mW (50.6 mA at 5 V). Energy per inference is
active power times measured latency. Values are directly measured with an INA219,
replacing the previous datasheet-based estimate.}
\label{tab:energy}
\begin{tabular}{llrrr}
\hline
Model & Dataset & Active (mW) & Latency (ms) & Energy (mJ) \\
\hline
TinyCNN8        & UT-HAR  & 354 & 14.1  & 5.0 \\
TinyMLP         & UT-HAR  & 301 & 33.3  & 10.0 \\
TinyCNN16       & UT-HAR  & 361 & 29.7  & 10.7 \\
TinyCNN32       & UT-HAR  & 363 & 67.8  & 24.6 \\
deepCNN         & UT-HAR  & 314 & 225.6 & 71.0 \\
deepTransformer & UT-HAR  & 347 & 169.8 & 58.9 \\
\hline
TinyCNN8        & CSI-HAR & 369 & 9.4   & 3.5 \\
TinyMLP         & CSI-HAR & 305 & 19.6  & 6.0 \\
TinyCNN16       & CSI-HAR & 381 & 20.2  & 7.7 \\
TinyCNN32       & CSI-HAR & 377 & 46.4  & 17.5 \\
deepCNN         & CSI-HAR & 331 & 243.5 & 80.7 \\
deepTransformer & CSI-HAR & 367 & 320.7 & 117.8 \\
\hline
\end{tabular}
\end{table}

% Factual note (author to reword): the measured active current (73-77 mA for the
% CNN and Transformer models) exceeds the 68 mA datasheet figure, which confirms
% that the earlier 50 mA assumption was not a conservative upper bound. The
% statistics MLP draws the least active power (about 300-305 mW).

% ---------------------------------------------------------------------------
% TABLE: cross-subject generalization gap across three datasets (Major #5).
% Random (subject-overlapping) vs subject-independent evaluation. The gap
% reproduces on datasets with far more subjects than the original three-user set.
% ---------------------------------------------------------------------------
\begin{table}[t]
\centering
\caption{Cross-subject generalization gap on three WiFi-CSI datasets. Moving from
a random (subject-overlapping) split to a subject-independent split reduces
accuracy by a similar large margin regardless of subject count, confirming that
the effect is not an artefact of the small three-user dataset.}
\label{tab:gengap-ext}
\begin{tabular}{lrrr}
\hline
Dataset (subjects) & Random (\%) & Subject-indep. (\%) & Gap (pt) \\
\hline
CSI-HAR (3)    & 92.1 & 62.1 & 30.0 \\
MM-Fi (10)     & 47.0 & 6.7  & 40.3 \\
MM-Fi (40)     & 35.8 & 4.8  & 31.0 \\
CSI-Bench (43) & 84.1 & 51.7 & 32.4 \\
\hline
\end{tabular}
\end{table}

% Factual note (author to reword): MM-Fi is a 27-class daily-activity HAR dataset;
% CSI-Bench is a 4-class motion-source task with 43 users. MM-Fi is reported both at
% 10 subjects (one environment) and at the full 40 subjects across all four
% environments (all 320,760 wifi-csi frames; a 456-dim per-subcarrier amplitude-stat
% RandomForest, 21,600 windows). In both cases the subject-independent accuracy falls
% to near the 27-class chance level (6.7\% at 10 subjects, 4.8\% at 40 subjects vs a
% 3.7\% chance), showing that apparently accurate models can fail almost completely on
% unseen users, and that the gap does not shrink as the subject pool grows to 40.
% Cite MM-Fi (Yang et al., CC BY-NC) and CSI-Bench.

% ---------------------------------------------------------------------------
% FIGURES
% ---------------------------------------------------------------------------
\begin{figure}[t]
\centering
\includegraphics[width=\columnwidth]{figs/fig_energy.png}
\caption{Measured energy per inference on the classic ESP32 (INA219, 240 MHz,
WiFi off) for all deployable models on both datasets.}
\label{fig:energy}
\end{figure}

\begin{figure}[t]
\centering
\includegraphics[width=\columnwidth]{figs/fig_gengap_external.png}
\caption{The cross-subject generalization gap reproduces across three datasets
spanning 3 to 43 subjects.}
\label{fig:gengap-ext}
\end{figure}

% ---------------------------------------------------------------------------
% TABLE: classic ESP32 vs ESP32-S3 (review #9). Same 12 int8 models, measured
% on both boards. S3 flashed via USB-C; energy = measured active power x measured
% latency. Place in a new results subsection "ESP32-S3 comparison".
% ---------------------------------------------------------------------------
\begin{table}[t]
\centering
\caption{Measured classic ESP32 vs ESP32-S3 on the same int8 models. Latency and
active power are measured on each board; energy is their product. The S3 (LX7 cores
with the esp-nn SIMD kernels and PSRAM) is on average 4.8x faster and uses 70\% less
energy per inference; the convolutional models benefit most from the vector kernels.}
\label{tab:s3}
\begin{tabular}{llrrrrr}
\hline
Model & Data & Classic (ms) & S3 (ms) & Speedup & Classic (mJ) & S3 (mJ) \\
\hline
TinyCNN8        & UT  & 14.1  & 2.4   & 5.9x & 5.0   & 0.66 \\
TinyMLP         & UT  & 33.3  & 25.3  & 1.3x & 10.0  & 6.81 \\
TinyCNN16       & UT  & 29.7  & 4.1   & 7.2x & 10.7  & 1.24 \\
TinyCNN32       & UT  & 67.8  & 8.7   & 7.7x & 24.6  & 2.88 \\
deepCNN         & UT  & 225.6 & 29.1  & 7.8x & 70.9  & 8.91 \\
deepTransformer & UT  & 169.8 & 115.5 & 1.5x & 58.9  & 28.85 \\
TinyCNN8        & CSI & 9.4   & 2.0   & 4.6x & 3.5   & 0.51 \\
TinyMLP         & CSI & 19.6  & 14.7  & 1.3x & 6.0   & 4.16 \\
TinyCNN16       & CSI & 20.2  & 3.5   & 5.8x & 7.7   & 0.98 \\
TinyCNN32       & CSI & 46.4  & 6.7   & 6.9x & 17.5  & 2.18 \\
deepCNN         & CSI & 243.5 & 40.2  & 6.1x & 80.7  & 11.86 \\
deepTransformer & CSI & 320.7 & 266.0 & 1.2x & 117.8 & 77.04 \\
\hline
\end{tabular}
\end{table}

\begin{figure}[t]
\centering
\includegraphics[width=\columnwidth]{figs/fig_classic_vs_s3.png}
\caption{Measured energy per inference, classic ESP32 vs ESP32-S3, with the latency
speedup labeled per model.}
\label{fig:s3}
\end{figure}

% ---------------------------------------------------------------------------
% TABLE: frequency and WiFi ablations (Kim "very important"). Active power measured
% at 160 vs 240 MHz, and with the WiFi stack enabled vs disabled (TinyCNN16).
% ---------------------------------------------------------------------------
\begin{table}[t]
\centering
\caption{Frequency and WiFi ablations on the classic ESP32 (CSI-HAR, measured active
power). Lowering the clock to 160\,MHz reduces active power by about 30\%, but raises
latency, so energy per inference is nearly unchanged (for TinyCNN16, 7.9\,mJ at
160\,MHz vs 7.7\,mJ at 240\,MHz); energy therefore does not scale purely with latency.
Bringing up the WiFi stack (station mode, idle) does not materially change the
per-inference active power.}
\label{tab:ablation}
\begin{tabular}{lrr}
\hline
Model & Active @160\,MHz (mW) & Active @240\,MHz (mW) \\
\hline
TinyCNN16       & 265 & 381 \\
TinyCNN32       & 266 & 377 \\
deepCNN         & 238 & 331 \\
deepTransformer & 264 & 367 \\
\hline
\multicolumn{3}{l}{WiFi ablation (TinyCNN16, 240\,MHz): off 381\,mW vs on 379\,mW} \\
\hline
\end{tabular}
\end{table}

% ---------------------------------------------------------------------------
% TABLE: deployment resource audit (Major #7 / spec #13). Compiled application
% flash vs model-file size vs tensor arena. Bootloader ~26 kB + partition ~3 kB fixed.
% ---------------------------------------------------------------------------
\begin{table}[t]
\centering
\caption{Deployment resource audit on the classic ESP32. Compiled application flash
(the full firmware image including the TensorFlow Lite Micro runtime and the model),
model-file size (int8 weights), and tensor arena (SRAM). Reporting compiled flash, not
just the model file, addresses the true on-device footprint; all fit well within the
4\,MB flash and the ~150\,kB SRAM arena budget.}
\label{tab:resource}
\begin{tabular}{llrrr}
\hline
Model & Data & Model file (kB) & App flash (kB) & Arena (kB) \\
\hline
TinyCNN8        & UT  & 11.1  & 434 & 12.9 \\
TinyMLP         & UT  & 184.0 & 606 & 12.5 \\
TinyCNN16       & UT  & 18.7  & 441 & 13.0 \\
TinyCNN32       & UT  & 37.6  & 460 & 13.4 \\
deepCNN         & UT  & 102.9 & 525 & 10.0 \\
deepTransformer & UT  & 88.3  & 510 & 24.0 \\
TinyCNN8        & CSI & 9.1   & 434 & 8.1 \\
TinyMLP         & CSI & 108.0 & 530 & 7.7 \\
TinyCNN16       & CSI & 14.6  & 437 & 8.3 \\
TinyCNN32       & CSI & 29.3  & 451 & 8.7 \\
deepCNN         & CSI & 86.3  & 508 & 12.0 \\
deepTransformer & CSI & 76.3  & 499 & 54.0 \\
\hline
\end{tabular}
\end{table}

% ---------------------------------------------------------------------------
% TABLE: firmware image audit (spec item F4 / #13). Segment sizes from idf.py size and the
% dominant per-archive contributors from idf.py size-components (classic ESP32, WiFi off).
% ---------------------------------------------------------------------------
\begin{table}[t]
\centering
\caption{Firmware image audit on the classic ESP32 (from the build tool). The TensorFlow Lite
Micro runtime is the single largest contributor to the image (about 232 kB of flash), far larger
than any int8 model file; static SRAM use is small (12 kB of .data plus .bss), leaving about 168 kB
of internal SRAM for the heap and tensor arena. This is why the deep models still fit despite the
absence of PSRAM.}
\label{tab:firmware-audit}
\begin{tabular}{lrr}
\hline
Segment / component & Flash (kB) & SRAM (kB) \\
\hline
Code .text (flash)        & 302.6 & --- \\
Const .rodata (incl. model) & 74.2 & --- \\
Static .data + .bss         & --- & 11.8 \\
IRAM .text                  & --- & 44.4 \\
\hline
esp-tflite-micro (runtime)  & 226.7 & --- \\
C library                   & 45.4  & --- \\
Application (this firmware)  & 21.6  & --- \\
FreeRTOS                    & 21.1  & --- \\
\hline
Total image                 & 430.3 & --- \\
SRAM free for heap + arena  & --- & 164.7 \\
\hline
\end{tabular}
\end{table}

% ---------------------------------------------------------------------------
% TABLE: tensor-arena RAM breakdown (spec item F3 / Major #7). Measured on-device
% with the RecordingMicroInterpreter (RECORD_ALLOC=1), which reports how the tensor
% arena splits across categories. Deterministic: depends only on the model graph and
% the TFLM version, and re-captured to identical bytes. Boot budget on the classic
% ESP32 (WiFi/BT off): 153.6 kB arena from a 155.6 kB largest free internal block.
% ---------------------------------------------------------------------------
\begin{table}[t]
\centering
\caption{On-device tensor-arena RAM breakdown for the deep models (classic ESP32,
measured with the recording interpreter). The arena splits into a non-persistent
head (peak activation/working memory, reused across layers) and a persistent tail
(operator scratch buffers, persistent tensors, quantization parameters, and the
per-operator bookkeeping structs). For the Transformer the activation head and the
operator scratch dominate, which is why its arena grows with the input window
(CSI-HAR 54 kB vs UT-HAR 24 kB); the CNN arena stays near 10-12 kB. All fit well
within the 153.6 kB arena budget.}
\label{tab:ram-breakdown}
\begin{tabular}{llrrrrr}
\hline
Model & Data & Arena (kB) & Head/act. (kB) & Tail/persist. (kB) & Op scratch (kB) & Ops \\
\hline
Deep CNN    & UT  & 10.1 & 6.0  & 4.1  & 2.6  & 16 \\
Deep CNN    & CSI & 12.1 & 8.0  & 4.1  & 2.6  & 16 \\
Transformer & UT  & 24.1 & 10.0 & 14.0 & 10.4 & 55 \\
Transformer & CSI & 54.1 & 36.0 & 18.0 & 14.4 & 55 \\
\hline
\end{tabular}
\end{table}

% Factual note (author to reword): the persistent tail also holds the small persistent
% tensors (80 B), quantization parameters (40 B), and NodeAndRegistration structs
% (512 B for the 16-op CNN, 1760 B for the 55-op Transformer); the remainder of the
% arena is the non-persistent activation head. Peak RAM is the arena_used figure, all
% far below the ~150 kB usable internal SRAM.

% ---------------------------------------------------------------------------
% TABLE: classical pipeline on-device timing (Major #7). The 5-statistic feature
% extraction dominates; tree traversal alone understates the true cost.
% ---------------------------------------------------------------------------
\begin{table}[t]
\centering
\caption{Classical tier on-device timing on the classic ESP32 (emlearn C). The
feature-extraction step (computing the five statistics over a $T{=}64$ window)
dominates the classical pipeline; Decision Tree and Random Forest traversal are
microseconds, so reporting traversal alone understates the true cost. The
statistics MLP predict is the slowest classical model.}
\label{tab:classical-timing}
\begin{tabular}{lrrrr}
\hline
Dataset & Feature extr. (ms) & DecisionTree ($\mu$s) & RandomForest ($\mu$s) & MLP-stats (ms) \\
\hline
UT-HAR  & 1.637 & 1.24 & 33.6 & --- \\
CSI-HAR & 0.885 & 1.17 & 7.5  & 4.67 \\
\hline
\end{tabular}
\end{table}

% ---------------------------------------------------------------------------
% TABLE: classical sklearn vs emlearn-C parity (spec item E2 / checklist #2). Exact-match
% rate of the emlearn C export against the trained scikit-learn model over the real test set,
% using the emlearn float export. Confirms the classical deployment reproduces the reference.
% ---------------------------------------------------------------------------
\begin{table}[t]
\centering
\caption{Output parity between the scikit-learn classical models and their emlearn C export,
measured as the exact-match rate over the full real test set (float export). The Decision Tree
and statistics MLP are bit-exact on both datasets; the Random Forest differs on 1.4\% of UT-HAR
windows only, because scikit-learn averages the per-tree probabilities before the argmax while the
C export resolves ties by tree vote. The compact int8 tree export is reserved for the latency
benchmark; the float export is used for accuracy-faithful deployment.}
\label{tab:emlearn-parity}
\begin{tabular}{llrr}
\hline
Dataset & Model & Exact match vs sklearn (\%) & Test windows \\
\hline
CSI-HAR & Decision Tree   & 100.0 & 420 \\
CSI-HAR & Random Forest   & 100.0 & 420 \\
CSI-HAR & MLP (statistics)& 100.0 & 420 \\
UT-HAR  & Decision Tree   & 100.0 & 500 \\
UT-HAR  & Random Forest   & 98.6  & 500 \\
UT-HAR  & MLP (statistics)& 100.0 & 500 \\
\hline
\end{tabular}
\end{table}

% ---------------------------------------------------------------------------
% TABLE: on-device vs desktop int8 numerical parity (spec item E3). The deployed
% TinyCNN16 (CSI-HAR) is run on the classic ESP32 (esp-tflite-micro int8 kernels)
% and on the desktop reference interpreter over the same seven stored real CSI
% windows; the full int8 output vector is compared element-wise. This verifies that
% the deployment does not silently diverge from the reference model.
% ---------------------------------------------------------------------------
\begin{table}[t]
\centering
\caption{On-device versus desktop int8 numerical parity for the deployed TinyCNN16
(CSI-HAR), over seven stored real CSI windows (one per class). The on-device
predictions match the desktop reference on all seven windows, and the full int8
logit vectors agree to within a single quantization step (1 LSB) on every element,
with three of seven windows bit-exact. The small residual comes from the optimized
on-device int8 kernels rounding differently from the reference kernels; it never
changes a prediction.}
\label{tab:parity}
\begin{tabular}{llcc}
\hline
Window & True class & Argmax match & Max logit diff (LSB) \\
\hline
0 & bend     & yes & 1 \\
1 & fall     & yes & 1 \\
2 & lie down & yes & 0 \\
3 & run      & yes & 0 \\
4 & sitdown  & yes & 0 \\
5 & standup  & yes & 1 \\
6 & walk     & yes & 1 \\
\hline
\multicolumn{4}{l}{Argmax agreement 7/7; bit-exact logits 3/7; worst-case logit diff 1 LSB.} \\
\hline
\end{tabular}
\end{table}
